\begin{figure}[ht]\centering
\begin{tikzpicture}[node distance=14mm and 18mm,>=Latex]
\node (pair1) [draw,rounded corners] {$I,\ J$};
\node (sum) [draw,rounded corners,right=of pair1] {$I+J$};
\node (geom1) [draw,rounded corners,right=of sum] {$X\cap Y$};
\node (pair2) [draw,rounded corners,below=of pair1] {$I,\ J$};
\node (cap) [draw,rounded corners,right=of pair2] {$I\cap J$};
\node (geom2) [draw,rounded corners,right=of cap] {$X\cup Y$};
\draw[->] (pair1)--node[above] {sum}(sum);\draw[->] (sum)--node[above] {defines}(geom1);
\draw[->] (pair2)--node[above] {intersection}(cap);\draw[->] (cap)--node[above] {defines}(geom2);
\end{tikzpicture}
\caption{Ideal operations reverse geometric behavior: sums define scheme-theoretic intersections, while intersections of ideals define unions.}\label{fig:vi37-03}\end{figure}
